\section{Discussion and conclusion}\label{sec:discussion}

\paragraph{What the value adds to the calculus.}
The set-valued theory asks whether an information state is viable; the safety value says how viable, and the two agree exactly where they should. On deterministic models the value-one level is the feedback viable-set recursion, the deficit is the lightest mass that any policy must abandon, the $\alpha$-vectors are the maximal survivable subsets, and the obstruction certificates are the deficit's lower-bound sides. The agreement is not merely qualitative: on the delayed hidden-regime class the certified table, census and witnesses are read off two closed forms, and the one place where the set-valued verdicts seemed to overreach---the timing certificate firing on cells where a time-varying schedule is safe---turns out to be a precise statement about the declared hold class. Two levers are separated and priced on the same instance: flexibility within a blind window and earlier observation each buy $\min(w,1-w)$, on disjoint cells. For a regulator the distinction is operational: on the timing cells a pre-committed alternating schedule is as good as early information, and on the common-action cells nothing but earlier information helps.

\paragraph{What the stochastic layer changes.}
With noisy sensing the deficit is a weighted sum of fatal-path probabilities, the minimal-mass bound is a statement about the deterministic limit, and the value becomes a design quantity for observation: a probe pinned to a deadline moves the viability threshold linearly in the deadline, a general schedule has an exact threshold which is additive exactly when the worst losses align, and repeated noisy readings purchase reliability at a fixed exchange rate in floor margin. These are small instances, but the laws are exact and the mechanisms---blind decline, probe excursion, wrong-action drift---are the ones that any larger model must contain.

\paragraph{Delimitations.}
The semantics is nominal expectation under a declared model; the robust quantifiers of the companion papers are recovered only at the deterministic degeneration, which is a change of operator, and the contamination family is where worst-case wording is given a formal home. Timing statements are class-scoped. No learning theory is claimed beyond the one-probe exact split and the majority count: Bayesian updating of continuous parameters, repeated experiments with design, and exploration--exploitation trade-offs are out of scope, and the hidden parameter enters additively. The $\alpha$-set growth is exponential in the worst case and the hardness results stand; the certified families with two or three witnesses are structural, not representative. The Sperner bound is cited, not formalized.

\paragraph{Outlook.}
Three directions follow directly. Exact point-based value iteration with rational solvers and survivable-set compression over the antichain would carry the deterministic layer to medium instances with several hidden parameters, where the Sperner bound rather than the action count governs witness growth; the computational companion \cite{Abaee2026cert} carries the first steps. The schedule law of Theorem~\ref{thm:additive} invites a treatment of misaligned worst losses and state-dependent parameter dynamics. And the contamination operator, exact at one stage, suggests ambiguity sets under which the multi-stage value remains computable in closed form on structured classes.

\paragraph{Conclusion.}
The belief-state safety value carries viability under incomplete observation from verdicts to numbers without loss of exactness. Its value-one level is the set-valued theory; below it, the deficit is a rational number with an explicit witness, the certificates are its lower bounds, the policy class is a priced design choice, and late or noisy information has an exact cost on the state. Every number in the paper is open to script-level re-inspection, and the structural theorems explain why the scripts return what they return.
